\appendixchapter{Quotient Spaces, Gluing, and Classical Manifolds}
\label{app:quotients}
\chaptermark{Quotient Spaces and Classical Manifolds}
\section*{Introduction}
This appendix makes the gluing pictures of Chapter~13 precise. It uses
Chapter~10's open-set, product, compactness, and Hausdorff language, and
Chapter~13's definitions of smooth manifolds and manifolds with boundary.
The analysis-to-Stokes route uses the product, subspace, and chart
constructions already given and does not depend on this appendix.
The construction begins by specifying which points represent the same point,
then determines continuity and local coordinates from that specification.
These examples prepare the global questions of geometric topology and
provide spaces on which the cohomological ideas of Appendix~\ref{app:cohomology}
can be tested. The quotient tools are structural constructions used repeatedly,
not additional hypotheses in the main integration theorems.

\section{Quotient Topology and Quotient Maps}\label{sec:quotient-topology}
Appendix~\ref{app:tensors} forms quotient vector spaces by declaring vectors equivalent.
Here we retain that idea but allow any equivalence relation on any topological
space. The extra question is which subsets of the resulting set should be open.
The progression is from quotient vector spaces to topological-group quotients
to arbitrary equivalence-relation quotients: increasing generality does not
mean that every quotient has a vector-space or group structure.

\begin{definition}[Quotient topology]
Let $\sim$ be an equivalence relation on a topological space $X$.
Its quotient set $X/{\sim}$ consists of the classes $[x]$.
The \textbf{canonical projection} and the \textbf{quotient topology} are
\[
\pi:X\longrightarrow X/{\sim},\quad \pi(x)=[x],\qquad
U\text{ open in }X/{\sim}\ \Longleftrightarrow\ \pi^{-1}(U)\text{ open in }X.
\]
Here $\pi^{-1}(U)$ denotes an inverse image, not an inverse function.
\end{definition}
\begin{proposition}[The open-set rule defines a topology]
The quotient topology is the finest topology making $\pi$ continuous.
\end{proposition}
\begin{proof}
Inverse images take the empty set and whole quotient to $\varnothing$ and
$X$, and preserve arbitrary unions and finite intersections. The open-set
axioms in $X$ therefore imply the three topology axioms on the quotient.
The rule makes $\pi$ continuous. Any topology making $\pi$ continuous can
contain only sets whose inverse images are open, so it is contained in this one.
\end{proof}
\begin{definition}[Quotient map]
A surjection $q:X\to Q$ between topological spaces is a \textbf{quotient map}
if $U\subseteq Q$ is open exactly when $q^{-1}(U)$ is open in $X$.
\end{definition}
We reserve $\pi$ for canonical projections onto equivalence classes and
$q$ for general quotient maps; $Q$ need not literally be written $X/{\sim}$.

\begin{proposition}[The universal property]\label{prop:quotient-universal}
For any topological space $Y$, a map $h:X/{\sim}\to Y$ is continuous
if and only if $h\circ\pi$ is continuous. Consequently, every continuous
$f:X\to Y$ constant on equivalence classes has a unique continuous factor
\[
\widetilde f:X/{\sim}\to Y,\qquad
\widetilde f([x])=f(x),\qquad f=\widetilde f\circ\pi.
\]
\end{proposition}
\begin{proof}
For open $V\subseteq Y$, $(h\circ\pi)^{-1}(V)=\pi^{-1}(h^{-1}(V))$.
The quotient rule makes this open exactly when $h^{-1}(V)$ is open.
For the factorization, constancy on classes makes the displayed rule
well defined; surjectivity of $\pi$ gives uniqueness, and the equivalence
just proved gives continuity.
\end{proof}
\begin{center}
\begin{tikzcd}[column sep=large,row sep=large]
X\arrow[r,"\pi"]\arrow[dr,"f"']&X/{\sim}\arrow[d,dashed,"\widetilde f"]\\
&Y
\end{tikzcd}
\end{center}
Construct $f$ before gluing and check that its values agree at glued points;
the dashed arrow is the uniquely induced map. We use a tilde for a map
produced by such a factorization, while retaining structural names such as
$T_1\otimes T_2$ and $\Lambda^kT$ in Appendix~\ref{app:tensors}.
The same inverse-image proof works for every quotient map.
The factor need not be injective: a constant map may collapse distinct
classes. It is injective exactly when $f(x)=f(y)$ implies $x\sim y$,
so that the classes are precisely the fibers of $f$.

\begin{example}[Coset quotients]
Let $G$ be an abelian group written additively, equipped with a topology
for which addition and negation are continuous, and let $H\leq G$ be a subgroup.
Then
\[
x\sim y\ \Longleftrightarrow\ x-y\in H,\qquad [x]=x+H,\qquad
\pi:G\to G/H.
\]
The subgroup laws give an equivalence relation; equip its classes with the
quotient topology above. This is a special case of the same definition.
If $V$ is a topological vector space and $W$ a linear subspace, it equips
Appendix~\ref{app:tensors}'s quotient vector space $V/W$ with a topology. For the torus,
$\Z^2$ is an additive subgroup of $\R^2$, but it is not an $\R$-linear
subspace (scalar multiplication by
\[
1/2
\]
 does not preserve it).
\end{example}

\Needspace{10\baselineskip}
\section{Quotient-Space Tools}\label{sec:quotient-tools}
\begin{proposition}[Compact quotients]
A quotient of a compact space is compact.
\end{proposition}
\begin{proof}
Its canonical projection is continuous and onto, so the continuous-image
compactness theorem applies.
\end{proof}
\paragraph{Separation can fail.} The quotient need not be Hausdorff even when the
original space is. For example, in $\R$ collapse all of $(0,1)$ to a single
class and leave other points separate. If $[0]$ and the collapsed class had
disjoint open neighborhoods, the inverse image of the first would be an
open set containing $0$ and hence meeting $(0,1)$, whereas the inverse image
of the second contains all of $(0,1)$. This is impossible.

Recall \cref{def:hausdorff}: in a Hausdorff space, any two distinct
points have disjoint open neighborhoods. Closed points alone do not
ensure this separation, as \cref{ex:zariski-line} shows. The following
criterion uses only the compactness results of Chapter~10.
\begin{proposition}[Compact-to-Hausdorff criterion]\label{prop:quotient-compact-bijection}
A continuous bijection $h:K\to H$ from a compact space to a Hausdorff
space is a homeomorphism. No Hausdorff hypothesis on $K$ is needed.
\end{proposition}
\begin{proof}
A closed subset $F$ of $K$ is compact: express a relatively open cover
of $F$ using open subsets of $K$, adjoin $K\setminus F$, use compactness
of $K$, and discard that extra set.
Its image $h(F)$ is compact and hence closed in $H$ by
\cref{thm:compact-closed}. Thus $h$ is a closed bijection, so
$(h^{-1})^{-1}(F)=h(F)$ is closed for every closed $F$. This proves
continuity of $h^{-1}$.
\end{proof}
For the examples below, use this five-step strategy:
\begin{enumerate}
\item Build a continuous map on the original space.
\item Check that it is constant on the prescribed equivalence classes.
\item Descend it by \cref{prop:quotient-universal}.
\item Check that the descended map is bijective, by identifying its fibers.
\item Prove the source compact and target Hausdorff, then apply
\cref{prop:quotient-compact-bijection}.
\end{enumerate}
Compactness may come from a compact set of representatives even when the
original space is not compact. This is the key to the next two examples.

\begin{proposition}[A metric for separated isometric identifications]
\label{prop:isometric-quotient-metric}
Let $Y\subseteq\R^d$ have its relative Euclidean metric. Let $\mathcal G$
be a family of bijective isometries of $Y$, containing the identity and
closed under composition and inverse. Suppose that for every $p,q\in Y$
and $R>0$ the set
\[
\{g\in\mathcal G:\|p-g(q)\|\leq R\}
\]
is finite. Identify $p$ and $q$ when $p=g(q)$ for some $g$. Then
\[
d([p],[q])=\min_{g\in\mathcal G}\|p-g(q)\|
\]
is a metric inducing the quotient topology. If in addition no nonidentity
map fixes a point, every point has a sufficiently small relative ball
that projects homeomorphically onto an open quotient neighborhood.
\end{proposition}
\begin{proof}
The idea is to find a nearest representative and then compare small balls.
Only finitely many maps can improve the identity's distance, so the minimum
exists. Replacing $p,q$ by $a(p),b(q)$ reindexes it by the bijection
$g\mapsto a^{-1}gb$ of $\mathcal G$. An attained zero means equal classes;
inverses give symmetry. If $g,h$ minimize for $(p,q)$ and $(q,r)$, then
\[
d([p],[r])\leq\|p-gh(r)\|
\leq\|p-g(q)\|+\|q-h(r)\|.
\]
This proves the triangle inequality.

The inverse image of a metric ball about $[p]$ is the union of the relative
balls of the same radius about $g(p)$, hence is open in $Y$. Conversely,
if $\pi^{-1}(U)$ is open and contains $p$, choose a relative ball
$B_Y(p,\rho)\subseteq\pi^{-1}(U)$. Every class at distance less than
$\rho$ has a representative in this ball and belongs to $U$.
Thus the two topologies agree. Also $\pi$ is open: for open $A\subseteq Y$,

\[
\pi^{-1}(\pi(A))=\bigcup_g g(A)
\]
 is open.

For the last assertion, finiteness with $p=q$ and $R=1$, together with
the absence of fixed points, gives $\delta_p>0$ such that
$\|p-g(p)\|\geq\delta_p$ for every nonidentity $g$.
If $0<2\rho<\delta_p$, the balls $B_Y(p,\rho)$ and
$g(B_Y(p,\rho))$ are disjoint: an intersection would put their centers
at distance less than $2\rho$. The projection is consequently injective
on this ball, and its continuity and openness make the restriction a
homeomorphism onto an open neighborhood.
\end{proof}

\begin{lemma}[Separated translates give quotient neighborhoods]
\label{lem:separated-quotient-neighborhood}
Let a family $\mathcal G$ of isometries of a Euclidean subspace $Y$
contain the identity and be closed under composition and inverse. Identify
points carried to each other by these maps, and give $Y/{\sim}$ its
quotient topology. If a relatively open ball $D$ satisfies
$D\cap g(D)=\varnothing$ for every nonidentity $g\in\mathcal G$, then
$\pi|_D:D\to\pi(D)$ is a homeomorphism and $\pi(D)$ is open.
\end{lemma}
\begin{proof}
If $x,y\in D$ have the same class, then $y=g(x)$ for some $g$.
Disjointness forces $g$ to be the identity, hence $x=y$. This proves
injectivity. For any relatively open $A\subseteq D$, the set $A$ is
open in $Y$ because $D$ is, and the saturation has the exact description
\[
\pi^{-1}(\pi(A))=\bigcup_{g\in\mathcal G}g(A).
\]
Both inclusions follow by writing equality of classes as $y=g(a)$.
Every $g(A)$ is open since an isometry with inverse is a homeomorphism.
Thus $\pi(A)$ is open by the quotient definition
(\cref{sec:quotient-topology}). The restriction is continuous and bijective
onto its image; its openness proves continuity of its inverse. Taking
$A=D$ proves that its image is an open neighborhood.
\end{proof}

\section{Circles and Spheres as Quotients}\label{sec:quotient-circle}
\paragraph{Construction and descent.}
For $s,t\in\R$, write $s\sim t$ if $s-t\in\Z$. The continuous map
\[
e(t)=(\cos2\pi t,\sin2\pi t)
\]
is onto $S^1$, and $e(s)=e(t)$ exactly when $s-t\in\Z$.
It therefore descends to a continuous bijection
$\widetilde e:\R/\Z\to S^1$.
\paragraph{Topology.}
Every class has a representative in $[0,1]$,
so the restriction of the canonical projection to $[0,1]$ is a continuous
surjection. Thus $\R/\Z$ is compact. The circle is Hausdorff as a
subspace of $\R^2$. The criterion proves
\[
\R/\Z\cong S^1.
\]
\paragraph{The finite interval model.}
Restricting $e$ to $[0,1]$ gives exactly the endpoint relation $0\sim1$,
with every other class a singleton. The same factorization and
compact-to-Hausdorff argument give

\[
[0,1]/(0\sim1)\cong S^1.
\]
 Joining endpoints is therefore a precise
topological construction, illustrated in Figure~\ref{fig:circle-gluing}.
The resulting circle has the local arc charts already used in Chapter~13;
the whole closed interval is not one chart.
\begin{figure}[htbp]
\centering
\begin{tikzpicture}[font=\small,>=Stealth]
\node at (1,1.5) {interval};
\draw[thick] (0,0)--(2,0);
\fill[purple!80!black] (0,0) circle(.065) node[below,text=black] {$0$};
\fill[purple!80!black] (2,0) circle(.065) node[below,text=black] {$1$};
\draw[->,thick] (2.5,0)--(3.2,0);
\begin{scope}[shift={(4.4,0)}]
\node at (0,1.5) {bend; join endpoints};
\draw[thick] (30:.85) arc(30:330:.85);
\foreach \a in {30,330}{\fill[purple!80!black] (\a:.85) circle(.065);}
\draw[->,purple!80!black] (.85,.4)--(.85,.08);
\draw[->,purple!80!black] (.85,-.4)--(.85,-.08);
\end{scope}
\draw[->,thick] (5.65,0)--(6.35,0);
\begin{scope}[shift={(7.6,0)}]
\node at (0,1.5) {circle};
\draw[thick] (0,0) circle(.85);
\fill[purple!80!black] (.85,0) circle(.065) node[right,text=black] {$[0]=[1]$};
\end{scope}
\end{tikzpicture}
\caption{Joining the interval's two purple endpoints gives one circle
point. A neighborhood of that point joins short intervals near both
endpoints. This explains the seam without mistaking the whole closed
interval for a single circle chart.}
\label{fig:circle-gluing}
\end{figure}

\FloatBarrier

\Needspace{13\baselineskip}
\subsection*{Collapsing the boundary of a disk}\label{sec:quotient-sphere}
Closing an interval joins its two endpoints. In higher dimensions we can
instead collapse the entire boundary to a single point.
\begin{proposition}[Sphere as a disk quotient]
For $n\geq1$, collapsing $\partial D^n$ to one class gives

\[
D^n/\partial D^n\cong S^n.
\]

\end{proposition}
\begin{proof}
Write $r=\|x\|$ and define the continuous map $f:D^n\to S^n$ by
\[
f(x)=\left(\frac{\sin(\pi r)}{r}x,\cos(\pi r)\right)\quad(r>0),
\qquad f(0)=(0,\ldots,0,1).
\]
Continuity at zero follows from
\[
\sin(\pi r)/r\to\pi.
\]

The whole boundary maps to the south pole. For $0<r<1$ the last
coordinate determines $r$ uniquely, and the first $n$ coordinates
determine the direction of $x$. To see surjectivity explicitly, write a
sphere point as $(u,t)$ with $\|u\|^2+t^2=1$. If $-1<t<1$, let

\[
r=\arccos(t)/\pi\in(0,1),\qquad x=ru/\sin(\pi r).
\]
 Since
$\|u\|=\sin(\pi r)$, this has $\|x\|=r$ and $f(x)=(u,t)$.
The north pole comes only from $x=0$; the south pole comes exactly from
the boundary. Thus the fibers are precisely the prescribed classes. The universal property gives
a continuous bijection
\[
\widetilde f:D^n/\partial D^n\to S^n.
\]

The domain is compact and the sphere Hausdorff, so it is a homeomorphism.
\end{proof}
Figure~\ref{fig:quotient-sphere} shows the two-dimensional case. This
construction also explains the sphere's one $0$-cell and one $n$-cell
in the optional CW discussion of \cref{sec:algebraic-topology-preview}.
\begin{figure}[htbp]
\centering
\begin{tikzpicture}[font=\small,>=Stealth]
\fill[gray!8] (0,0) circle(1);
\draw[blue!75!black,very thick] (0,0) circle(1);
\foreach \a in {0,45,...,315}{\fill[blue!75!black] ({cos(\a)},{sin(\a)}) circle(.045);}
\node[below] at (0,-1.15) {$D^2$: the whole boundary};
\draw[->,thick] (1.45,0)--(4.15,0) node[midway,above,align=center] {collapse every\\boundary point};
\draw[fill=gray!8] (5.7,0) circle(1.05);
\draw[dashed,gray] (4.65,0) arc[start angle=180,end angle=0,x radius=1.05,y radius=.3];
\draw[gray] (4.65,0) arc[start angle=180,end angle=360,x radius=1.05,y radius=.3];
\fill[blue!75!black] (5.7,-1.05) circle(.07);
\draw[->,blue!75!black] (6.8,-.8)--(5.8,-1.05);
\node[right] at (6.8,-.8) {one point};
\node[below] at (5.7,-1.35) {$S^2$};
\end{tikzpicture}
\caption{The rule for $D^2/\partial D^2$: every marked boundary point belongs
to one class, shown as the south pole; interior points remain distinct.}
\label{fig:quotient-sphere}
\end{figure}

\FloatBarrier

\Needspace{12\baselineskip}
\section{The Torus as a Quotient}\label{sec:quotient-torus}
The formal smooth torus of Chapter~13 is $T^2=S^1\times S^1$.
\paragraph{Construction and fibers.}
To identify the topology of the gluing model, define
\[
E(x,y)=\big((\cos2\pi x,\sin2\pi x),
                 (\cos2\pi y,\sin2\pi y)\big).
\]
This map is continuous and onto. Equality of its first circle coordinates
means the first parameters differ by an integer; equality of the second
means the same for the second parameters. Its fibers are therefore exactly
\[
(x,y)\sim(x+m,y+n),\qquad (m,n)\in\Z^2.
\]
The universal property gives a continuous bijection
$\widetilde E:\R^2/\Z^2\to S^1\times S^1$.
\paragraph{Topology.}
Every class has a representative in $[0,1]^2$, so the continuous image of
this compact square is the whole quotient. The target is Hausdorff:
if two pairs differ, disjoint neighborhoods in a differing factor,
together with the entire other factor, separate them. Thus
\cref{prop:quotient-compact-bijection} proves
\[
\R^2/\Z^2\cong S^1\times S^1.
\]
\paragraph{The square model.}
On the square, the fibers of $E$ are precisely the equivalence relation
generated by $(0,y)\sim(1,y)$ and $(x,0)\sim(x,1)$.
In particular the four corners form one class. Descending the restricted
map gives a continuous bijection from the compact square quotient to the
same Hausdorff product, hence a homeomorphism. This formally justifies
Figure~\ref{fig:torus-parameter-square}, including both stages of gluing.

\begin{figure}[htbp]
\centering
\begin{tikzpicture}[font=\small,>=Stealth]
\fill[gray!7] (0,0) rectangle(2,2);
\draw[red!75!black,very thick,->] (0,0)--(0,2);
\draw[red!75!black,very thick,->] (2,0)--(2,2);
\draw[blue!75!black,very thick,dashed,->] (0,0)--(2,0);
\draw[blue!75!black,very thick,dashed,->] (0,2)--(2,2);
\foreach \x/\y in {0/0,0/2,2/0,2/2}{\fill (\x,\y) circle(.05);}
\node[left] at (0,1) {$b$};\node[right] at (2,1) {$b$};
\node[below] at (1,0) {$a$};\node[above] at (1,2) {$a$};
\draw[->,thick] (2.8,1)--(4.3,1) node[midway,above] {identify};
\begin{scope}[shift={(6.3,1)},x={(1cm,0cm)},y={(0cm,.4cm)},z={(0cm,.75cm)}]
\foreach \v in {0,45,...,315}{\draw[gray!60] plot[domain=0:360,samples=65] ({(1.35+.42*cos(\v))*cos(\x)},{(1.35+.42*cos(\v))*sin(\x)},{.42*sin(\v)});}
\foreach \u in {0,30,...,330}{\draw[gray!45] plot[domain=0:360,samples=40] ({(1.35+.42*cos(\x))*cos(\u)},{(1.35+.42*cos(\x))*sin(\u)},{.42*sin(\x)});}
\end{scope}
\node[below] at (6.3,-.2) {$S^1\times S^1$};
\end{tikzpicture}
\caption{Torus quotient: $(0,y)\sim(1,y)$ and $(x,0)\sim(x,1)$.
The dashed $a$-edges and solid $b$-edges distinguish the two pairs.
The arrows agree on each paired edge, and all four vertices form one class.
Figure~\ref{fig:torus-parameter-square} gives the full geometric sequence.}
\label{fig:quotient-torus}
\end{figure}

\FloatBarrier

\paragraph{Coordinate conventions.}
The linear rescaling
\[
(\theta,\varphi)\mapsto
(\theta/(2\pi),\varphi/(2\pi))
\]
 and its inverse respect the respective
relations; both descend continuously by the universal property. Hence
$\R^2/(2\pi\Z)^2\cong\R^2/\Z^2$.
These homeomorphisms identify the topological models. The smooth structure
is the product-circle structure already constructed in Chapter~13,
transported through the homeomorphism if the quotient notation is used.
No general smooth-quotient theorem is needed.

\Needspace{13\baselineskip}
\section{Cylinder and M\"obius Strip}\label{sec:quotient-bands}
Gluing two edges leaves boundary; reversing that gluing changes the way
in which the remaining edges join.
\begin{definition}[Finite rectangle models]
For $\sigma\in\{1,-1\}$, put
\[
Q_\sigma=([0,1]\times[-1,1])/{\sim},\qquad (0,t)\sim(1,\sigma t),
\]
with only the equivalence relation generated by these edge identifications.
The \textbf{cylinder} is $Q_1$ and the \textbf{M\"obius strip} is $Q_{-1}$.
Both have the quotient topology and are compact.
\end{definition}
\begin{figure}[htbp]
\centering
\begin{tikzpicture}[font=\small,>=Stealth]
\fill[gray!7] (0,0) rectangle(2.3,1.5);
\draw[purple!80!black,very thick] (0,0)--(2.3,0) (0,1.5)--(2.3,1.5);
\draw[red!75!black,very thick,->] (0,0)--(0,1.5);
\draw[red!75!black,very thick,->] (2.3,1.5)--(2.3,0);
\node[left] at (0,.75) {$t$};\node[right] at (2.3,.75) {$-t$};
\node[below] at (1.15,0) {$0\leq x\leq1$};
\draw[->,thick] (3.1,.75)--(4.3,.75);
\begin{scope}[shift={(6.4,.75)},x={(1cm,0cm)},y={(.3cm,.4cm)},z={(0cm,.9cm)}]
\foreach \u in {0,12,...,348}{\draw[gray!50] ({(1.3-.38*cos(\u/2))*cos(\u)},{(1.3-.38*cos(\u/2))*sin(\u)},{-.38*sin(\u/2)})--({(1.3+.38*cos(\u/2))*cos(\u)},{(1.3+.38*cos(\u/2))*sin(\u)},{.38*sin(\u/2)});}
\foreach \a/\b/\sty in {0/180/solid,180/360/dashed,360/540/solid,540/720/dashed}{
\draw[purple!80!black,very thick,\sty] plot[domain=\a:\b,samples=50] ({(1.3+.38*cos(\x/2))*cos(\x)},{(1.3+.38*cos(\x/2))*sin(\x)},{.38*sin(\x/2)});}
\draw[red!75!black,very thick,->] (.92,0,0)--(1.68,0,0);
\end{scope}
\node[below] at (6.4,-.45) {one boundary circle};
\end{tikzpicture}
\caption{M\"obius quotient: $(0,t)\sim(1,-t)$. Opposite arrows encode the
reversal. The two long edges join into the single emphasized boundary;
dashes indicate its rear portions. Figure~\ref{fig:mobius-gluing} shows
the deformation, while Figure~\ref{fig:quotient-seam} checks local charts.}
\label{fig:quotient-mobius}
\end{figure}

\FloatBarrier

For the cylinder, $(x,t)\mapsto(e(x),t)$ has exactly these fibers. The
compact-to-Hausdorff criterion therefore identifies $Q_1$ with
$S^1\times[-1,1]$.

\subsection{An infinite strip for local coordinates}
The finite rectangle shows the global gluing; an infinite strip lets a
neighborhood continue across a seam without cutting its coordinates.
\begin{construction}[Infinite strip model]
Set $B=\R\times[-1,1]$ and define
\[
f_k(x,t)=(x+k,\sigma^k t),\qquad k\in\Z.
\]
These Euclidean isometries satisfy $f_kf_l=f_{k+l}$ and
$f_k^{-1}=f_{-k}$. Declare $p\sim r$ when $p=f_k(r)$ for some $k$;
the identities verify equivalence. Write $\widehat Q_\sigma=B/{\sim}$
and $\pi:B\to\widehat Q_\sigma$ for its canonical projection.
\end{construction}
To compare the two models by compact-to-Hausdorff, we first establish
Hausdorffness of the infinite model directly.
\begin{lemma}[Distance between strip classes]\label{lem:band-metric}
The formula
\[
d([p],[r])=\min_{k\in\Z}\|p-f_k(r)\|
\]
defines a metric on $\widehat Q_\sigma$.
\end{lemma}
\begin{proof}
A bound on $\|p-f_k(r)\|$ bounds the horizontal displacement and hence
allows only finitely many integers $k$. No $f_k$ with $k\ne0$ fixes a
point, and $\|p-f_k(p)\|\geq |k|\geq1$. Thus
\cref{prop:isometric-quotient-metric} applies. It also shows that this
metric induces the quotient topology and that relative balls of radius
less than
\[
1/2
\]
 project to quotient neighborhoods. In particular the
quotient is Hausdorff.
\end{proof}
\begin{proposition}[The finite and infinite models agree]
The rectangle inclusion induces a homeomorphism
$Q_\sigma\cong\widehat Q_\sigma$.
\end{proposition}
\begin{proof}
\emph{Construct the map.} Compose inclusion with $\pi$. It respects the
edge relation because $f_1(0,t)=(1,\sigma t)$, so
$\pi(0,t)=\pi(1,\sigma t)$. Constancy on these generating pairs implies
constancy on every equivalence class. The quotient universal property
therefore gives the continuous map $[x,t]\mapsto\pi(x,t)$.

\emph{Check the fibers.} Reducing the first coordinate modulo integers gives
a representative in the rectangle: for $x\in\R$ choose integer $k$
with $x-k\in[0,1]$ and apply $f_{-k}$. If rectangle points satisfy
$(x,t)=f_k(y,u)$, then $x-y=k$ with $x,y\in[0,1]$. Thus $k=0$ gives
identical points, while $k=\pm1$ forces the paired endpoints $0,1$ and
the specified sign change in the second coordinate. There are no other
possibilities. These are exactly the finite-model fibers, proving bijectivity.

\emph{Compare the topologies.} Its source is compact and its target is
Hausdorff by \cref{lem:band-metric}. Apply
\cref{prop:quotient-compact-bijection}.
\end{proof}

\subsection{Charts across seams and at boundary points}
\begin{proposition}[Small strip neighborhoods are quotient charts]
For $p\in B$ and
\[
0<\rho<1/4,
\]
 let
$D=B\cap B_{\R^2}(p,\rho)$. Then
$\pi|_D:D\to\pi(D)$ is a homeomorphism onto an open quotient neighborhood.
\end{proposition}
\begin{proof}
For $k\ne0$, the centers $p$ and $f_k(p)$ have horizontal separation
at least $1$. If $z\in D\cap f_k(D)$, then
\[
1\leq\|p-f_k(p)\|\leq\|p-z\|+\|z-f_k(p)\|<2\rho<1,
\]
a contradiction. Apply \cref{lem:separated-quotient-neighborhood} to the
relative ball and the maps $f_k$.
\end{proof}
At an interior point $p=(x_0,t_0)$ choose also $\rho<1-|t_0|$; $D$ is
an ordinary disk. At $t_0=1$ use $(x-x_0,1-t)$, and at $t_0=-1$ use
$(x-x_0,t+1)$; the image is an open half-disk relative to $v\geq0$.
These points give genuine boundary charts.

For a seam chart centered at $(0,t_0)$, let $D$ be such a relative disk.
In the finite rectangle use precisely the pieces
\[
D_0=D\cap\{x\geq0\},\qquad D_1=f_1(D\cap\{x\leq0\}).
\]
The coordinate map into $D$ is
\[
\chi([x,t])=
\begin{cases}(x,t),&(x,t)\in D_0,\\
(x-1,\sigma t),&(x,t)\in D_1.
\end{cases}
\]
At the identified edges $(0,t)\sim(1,\sigma t)$ both formulas give
$(0,t)$. Thus the map is well defined; the small-disk proposition proves
it is a homeomorphism. Interior seams assemble two half-disks into a disk;
seam endpoints assemble two quarter-disks into a half-disk.
\begin{figure}[htbp]
\centering
\begin{tikzpicture}[font=\small,>=Stealth]
\node at (1.6,2.65) {Finite rectangle: $\sigma=-1$};
\fill[gray!5] (0,0) rectangle(3.2,2.2);
\draw (0,0) rectangle(3.2,2.2);
\fill[blue!20] (0,1.55)--(0,2) arc(90:-90:.45)--cycle;
\fill[orange!30] (3.2,.65)--(3.2,.2) arc(-90:-270:.45)--cycle;
\draw[blue!75!black,thick,->] (0,1.15)--(0,1.95);
\draw[orange!85!black,thick,->] (3.2,1.05)--(3.2,.25);
\fill (0,1.55) circle(.04);\fill (3.2,.65) circle(.04);
\node[left] at (0,1.55) {$t_0$};\node[right] at (3.2,.65) {$-t_0$};
\node[below] at (0,0) {$x=0$};\node[below] at (3.2,0) {$x=1$};
\draw[->,thick] (4.2,1.1)--(5.55,1.1);
\node[align=center] at (4.85,2.05) {reflect $t$\\and translate};
\begin{scope}[shift={(7,1.1)}]
\fill[blue!20] (0,0)--(0,1) arc(90:-90:1)--cycle;
\fill[orange!30] (0,0)--(0,-1) arc(-90:-270:1)--cycle;
\draw[dashed] (0,0) circle(1);
\draw[densely dotted] (0,-1)--(0,1);
\draw[->,blue!75!black,thick] (.18,-.4)--(.18,.4);
\draw[->,orange!85!black,thick] (-.18,-.4)--(-.18,.4);
\fill (0,0) circle(.045);
\node at (-.6,0) {$D_1$};\node at (.6,0) {$D_0$};
\end{scope}
\node at (7,2.65) {One open coordinate disk};
\node[below,align=center] at (7,-.25) {$D_0:(x,t)$\\$D_1:(x-1,-t)$};
\end{tikzpicture}
\caption{A M\"obius seam chart. The right-edge piece is centered at
$(1,-t_0)$ and is reflected transversely before being placed to the left of
the piece at $(0,t_0)$. Both arrows then point in the same chart direction.
For the cylinder use $(x-1,t)$ instead, without reflection. At a boundary
endpoint the corresponding quarter-disks assemble into a half-disk.
Dashed arcs indicate excluded disk boundaries.}
\label{fig:quotient-seam}
\end{figure}

\FloatBarrier

\paragraph{Smoothness and second countability.}
On an overlap, lifts differ locally by one $f_k$; the integer is locally
constant: disjoint translated small balls make the chosen lift unique
on a neighborhood. It is constant on each connected overlap component
by \cref{prop:clopen-locally-constant,def:connected-component}. Transitions, including the boundary coordinates, are locally
affine and therefore smooth. Their maximal compatible atlas is the smooth
structure. The projection is open by the saturation argument above.
Second countability means having a countable topological basis
(\cref{def:second-countable}). By
\cref{ex:rational-basis,prop:second-countable-inheritance}, $B$ has a
countable relative Euclidean basis. Its images form a countable
basis: lift a point of an open quotient set and choose a basis element
inside its inverse image. Both spaces are consequently Hausdorff,
second countable smooth surfaces with boundary.

\begin{proposition}[Boundary components]
\[
\partial Q_1\cong S^1\sqcup S^1,\qquad \partial Q_{-1}\cong S^1.
\]
\end{proposition}
\begin{proof}
The charts identify the boundary as the image of $t=\pm1$.
For $\sigma=1$ the two levels remain separate; on each, parameters are
identified modulo $\Z$, giving two circles. For $\sigma=-1$, every boundary
class has a representative $(x,1)$, and two such representatives agree
exactly when their parameters differ by an even integer. Thus the map
$x\mapsto[x,1]$ descends to a continuous bijection
$\R/(2\Z)\to\partial Q_{-1}$. Its source is compact and its target is
Hausdorff, so it is a homeomorphism. Rescaling $x$ by
\[
1/2
\]
 identifies the
source with $\R/\Z\cong S^1$.
\end{proof}
Following the top rectangle edge continues along the bottom edge before
returning to the start, as traced in Figure~\ref{fig:mobius-gluing}.

\Needspace{12\baselineskip}
\section{The Klein Bottle}\label{sec:quotient-klein}
\begin{definition}[Square model]
With exactly the convention of Figure~\ref{fig:klein-gluing}, set
\[
K=[0,1]^2/{\sim},\qquad(0,y)\sim(1,y),\qquad(x,0)\sim(1-x,1).
\]
The equivalence relation is generated by these edge identifications.
All four corners form one class, and $K$ is compact.
\end{definition}
\begin{figure}[htbp]
\centering
\begin{tikzpicture}[font=\small,>=Stealth]
\fill[gray!7] (0,0) rectangle(2,2);
\draw[red!75!black,very thick,->] (0,0)--(0,2);
\draw[red!75!black,very thick,->] (2,0)--(2,2);
\draw[blue!75!black,very thick,->] (0,0)--(2,0);
\draw[blue!75!black,very thick,->] (2,2)--(0,2);
\node[below] at (1,0) {$x$};\node[left] at (0,1) {$y$};
\draw[->,thick] (2.7,1)--(4,1);
\begin{scope}[shift={(5.6,.2)},scale=.8]
% Shared bottle representation; the caller supplies scale and seam markers.
\fill[gray!8] (-.45,1.85) .. controls (-.5,1.1) and (-1,.9) .. (-1,.25)
 .. controls (-1,-.6) and (1,-.6) .. (1,.25)
 .. controls (1,.9) and (.45,1.1) .. (.45,1.85) --cycle;
\draw[thick] (-.45,1.85) .. controls (-.5,1.1) and (-1,.9) .. (-1,.25)
 .. controls (-1,-.6) and (1,-.6) .. (1,.25)
 .. controls (1,.9) and (.45,1.1) .. (.45,1.85);
\draw[thick] (-.45,1.85) .. controls (-.45,3.05) and (2.1,2.9) .. (2.1,1.5)
 .. controls (2.1,.9) and (.5,.8) .. (.15,.15);
\draw[thick] (.45,1.85) .. controls (.45,2.35) and (1.5,2.35) .. (1.5,1.5)
 .. controls (1.5,1.15) and (.25,1) .. (-.15,.15);

\draw[blue!75!black,very thick] (0,.12) ellipse(.2 and .1);
\node[below,blue!75!black] at (0,-.45) {glued seam};
\draw[orange!80!black,dotted,thick] (.85,.9) circle(.4);
\draw[->] (1.2,.9)--(2.35,.55);
\node[right,align=left] at (2.35,.55) {crossing only:\\no gluing};
\end{scope}
\end{tikzpicture}
\caption{Klein quotient: $(0,y)\sim(1,y)$ and $(x,0)\sim(1-x,1)$.
The solid blue seam represents the joined circles; the dotted callout
marks an incidental crossing of distinct sheets. The drawing in $\R^3$
is not an embedding. See Figure~\ref{fig:klein-gluing} for the deformation
and Figure~\ref{fig:klein-corner} for the local disk at the corner class.}
\label{fig:quotient-klein}
\end{figure}

\FloatBarrier

\subsection{The plane model and its topology}
As for the strip, extending coordinates past the seams makes neighborhoods
easier to construct. On $\R^2$ introduce the isometries
\[
G_{m,n}(x,y)=((-1)^n x+m,y+n),\qquad m,n\in\Z.
\]
They satisfy
\[
G_{m,n}G_{a,b}=G_{m+(-1)^na,n+b},\qquad
G_{m,n}^{-1}=G_{-(-1)^nm,-n}.
\]
Their orbits are equivalence classes; denote the quotient by $\widehat K$
and its projection by $\pi:\R^2\to\widehat K$.

\begin{proposition}[Quotient metric and Hausdorffness]
The formula
\[
d([p],[r])=\min_{m,n\in\Z}\|p-G_{m,n}(r)\|
\]
is a metric inducing the quotient topology on $\widehat K$.
The square inclusion induces a homeomorphism $K\cong\widehat K$.
\end{proposition}
\begin{proof}
\emph{Apply the common metric criterion.} A distance bound first bounds
the vertical integer $n$ and then the horizontal integer $m$, so only
finitely many pairs occur. For a nonidentity $G_{m,n}$, displacement is
at least $1$: use the vertical coordinate if $n\ne0$, and the horizontal
coordinate otherwise. The family therefore satisfies
\cref{prop:isometric-quotient-metric}, proving the metric and topology
claims and providing small quotient disks. No metric axioms need to be
reconstructed for this model.

\emph{Finite model.} Reduce $y$ and then $x$ to place any representative in
the square. In the square the only duplications are the prescribed edges
and corners: $G_{1,0}$ identifies vertical edges and $G_{1,1}$ identifies
$(x,0)$ with $(1-x,1)$. Inclusion therefore descends to a continuous
bijection from compact $K$ onto Hausdorff $\widehat K$, a homeomorphism.
\end{proof}

\Needspace{12\baselineskip}
\subsection{Interior, edge, and corner neighborhoods}
\begin{proposition}[Local Euclidean structure]
Every point of $K$ has an ordinary open-disk neighborhood, including the
single class represented by all four square corners.
\end{proposition}
\begin{proof}
For $(m,n)\ne(0,0)$ the displacement of a point under $G_{m,n}$ is at
least $1$: use the vertical displacement when $n\ne0$ and the horizontal
one otherwise. A disk of radius
\[
\rho<1/4
\]
 therefore has disjoint translates.
The triangle estimate in the strip proof shows the disk misses every
nontrivial translate. By \cref{lem:separated-quotient-neighborhood}, its
projection is a homeomorphism onto an open neighborhood.

\emph{Interior.} A sufficiently small disk remains inside the square and
already supplies the desired neighborhood.

\emph{Edges.} A disk centered on an edge away from the corners has two
half-disk pieces, represented at the paired edges. The edge rule fits
them into the full disk, not a half-disk.

\emph{Corner class.} Center the plane disk at $(0,0)$. Write its coordinates
as $(u,v)$. Its four quadrant pieces have square representatives
\[
\begin{array}{c|c}
\text{quadrant}&\text{representative in the square}\\\hline
u\geq0,\ v\geq0&(u,v)\\
u\leq0,\ v\geq0&(1+u,v)\\
u\geq0,\ v\leq0&(1-u,1+v)\\
u\leq0,\ v\leq0&(-u,1+v)
\end{array}
\]
On the axes these choices agree as classes by the edge relations.
Thus four quarter-disks assemble into one disk. There is no boundary
or corner singularity.
\end{proof}
\begin{figure}[htbp]
\centering
\begin{tikzpicture}[font=\small,>=Stealth]
\node at (1.35,3.25) {Four square-corner pieces};
\draw (0,0) rectangle(2.7,2.7);
\fill[blue!25] (0,0)--(.6,0) arc(0:90:.6)--cycle;
\fill[orange!35] (2.7,0)--(2.7,.6) arc(90:180:.6)--cycle;
\fill[green!25!white] (2.7,2.7)--(2.1,2.7) arc(180:270:.6)--cycle;
\fill[purple!25] (0,2.7)--(0,2.1) arc(270:360:.6)--cycle;
\node at (.22,.22) {A};\node at (2.48,.22) {B};
\node at (2.48,2.48) {C};\node at (.22,2.48) {D};
\draw[red!70!black,thick,->] (0,.8)--(0,1.8);
\draw[red!70!black,thick,->] (2.7,.8)--(2.7,1.8);
\draw[blue!75!black,thick,->] (.8,0)--(1.8,0);
\draw[blue!75!black,thick,->] (1.8,2.7)--(.8,2.7);
\foreach \p in {(0,0),(2.7,0),(0,2.7),(2.7,2.7)}{\fill \p circle(.035);}
\draw[->,thick] (3.3,1.35)--(4.55,1.35);
\node[align=center] at (3.9,2.2) {use the\\edge rules};
\begin{scope}[shift={(6,1.35)}]
\fill[blue!25] (0,0)--(1.1,0) arc(0:90:1.1)--cycle;
\fill[orange!35] (0,0)--(0,1.1) arc(90:180:1.1)--cycle;
\fill[purple!25] (0,0)--(-1.1,0) arc(180:270:1.1)--cycle;
\fill[green!25!white] (0,0)--(0,-1.1) arc(270:360:1.1)--cycle;
\draw[dashed] (0,0) circle(1.1);
\draw[dotted] (-1.1,0)--(1.1,0) (0,-1.1)--(0,1.1);
\node at (.45,.45) {A};\node at (-.45,.45) {B};
\node at (.45,-.45) {C};\node at (-.45,-.45) {D};
\fill (0,0) circle(.05);
\end{scope}
\node at (6,3.25) {One disk at the corner class};
\node[below,align=center] at (6,-.1) {all four corner marks\\become the center};
\end{tikzpicture}
\caption{The Klein-bottle corner chart. Bottom-left A stays in the positive
quadrant; bottom-right B translates left. Top-right C and top-left D reflect
horizontally and translate down, as in the displayed representative table.
The four pieces fill a disk with no missing quadrant, boundary, or singular
corner. The outer dashed circle is excluded.}
\label{fig:klein-corner}
\end{figure}

\FloatBarrier

\paragraph{Smooth structure.}
Lifted charts have transitions locally given by one $G_{m,n}$. These are
affine maps and hence smooth; their maximal extension is a smooth atlas.
\paragraph{Second countability.}
Projection is open since saturations are unions of isometric translates.
Images of the rational disks of \cref{ex:rational-basis} form a countable quotient basis,
by the same lift-and-shrink argument as for the strip. This checks the
remaining manifold hypothesis independently of the pictures.
\paragraph{The three-dimensional picture.}
The familiar bottle-like drawing in $\R^3$ self-intersects and is not an
embedding. Its crossing represents distinct sheets, not an additional
identification in $K$. The abstract compact smooth surface has the ordinary
disk neighborhoods just proved.


\section{Real Projective Space}\label{sec:projective-space}
To record a direction without choosing a length or a sign, replace each
nonzero vector by the line it spans. Ratios of coordinates will then give
ordinary local coordinates on the resulting space.
The proof architecture selectively follows the supplied smooth-manifolds
notes~\cite{xuewang-manifold-notes}, with explicit quotient factorizations
and chart inverses in the notation used here.
\begin{definition}[Scaling quotient and homogeneous coordinates]
For $n\geq0$, define
\[
\mathbb{RP}^n=(\R^{n+1}\setminus\{0\})/{\sim},\qquad
x\sim y\ \Longleftrightarrow\ y=\lambda x\text{ for some }\lambda\ne0,
\]
with the quotient topology and canonical projection
$\pi:\R^{n+1}\setminus\{0\}\to\mathbb{RP}^n$.
Write $[x_0:\cdots:x_n]$ for the class of $x$.
\end{definition}
The scalars $1$, $\lambda^{-1}$, and products of nonzero scalars verify
reflexivity, symmetry, and transitivity. Each class consists of the nonzero
vectors in exactly one one-dimensional linear subspace of $\R^{n+1}$.
Thus a projective point is a line through the origin.
Homogeneous coordinates are not ordinary global coordinates: common
nonzero rescaling describes the same projective point.

\subsection{The antipodal sphere model}
Every line meets $S^n$ in exactly two antipodal points. To identify the
topologies, rather than just the sets, we use the universal property twice.
\begin{proposition}[Scaling and antipodal quotients]
There is a natural homeomorphism
\[
\mathbb{RP}^n\cong S^n/(x\sim-x).
\]

\end{proposition}
\begin{proof}
Let
\[
\pi_S:S^n\to S^n/{\pm1}
\]
 be the antipodal projection and let

\[
N(x)=x/\|x\|.
\]
 The continuous composite $\pi_S\circ N$ is unchanged
by nonzero scaling, including negative scaling, so it induces a continuous
map
\[
\widetilde N:\mathbb{RP}^n\to S^n/{\pm1}.
\]

The inclusion $\iota:S^n\to\R^{n+1}\setminus\{0\}$ followed by $\pi$
is constant on antipodal pairs, so it induces a continuous map

\[
\widetilde\iota:S^n/{\pm1}\to\mathbb{RP}^n.
\]

Normalizing a unit vector changes nothing, and normalizing a nonzero
vector leaves its line unchanged. Thus both composites are identities.
\end{proof}
\begin{figure}[htbp]
\centering
\begin{tikzpicture}[font=\small,>=Stealth]
\draw[fill=gray!6] (0,0) circle(1.4);
\draw[dashed,gray] (-1.4,0) arc[start angle=180,end angle=0,x radius=1.4,y radius=.4];
\draw[gray] (-1.4,0) arc[start angle=180,end angle=360,x radius=1.4,y radius=.4];
\draw[thick,<->] (-1.7,-1.275)--(1.7,1.275);
\fill (0,0) circle(.04) node[below right] {$0$};
\foreach \x/\y/\lab/\pos in {1.12/.84/x/above left,-1.12/-.84/{-x}/below right}{
\fill[blue!75!black] (\x,\y) circle(.075);\node[\pos] at (\x,\y) {$\lab$};}
\node[left] at (-1.4,.65) {$S^2$};
\node[above right] at (1.7,1.275) {line $\R x$};
\draw[->,thick] (2.7,0)--(4.2,0) node[midway,above] {$\pi_S$};
\fill[blue!75!black] (4.8,0) circle(.075);
\node[right,align=left] at (5,0) {one point $[x]=[-x]$\\of $\mathbb{RP}^2$};
\node[align=center] at (2.5,-2) {a line through the origin $=$ an antipodal pair $=$ a projective point};
\end{tikzpicture}
\caption{The sphere model of projective space pairs $x$ only with $-x$.
The matching filled markers lie on the same line through the origin.
The final dot represents one class, not a drawing of the whole quotient.}
\label{fig:projective-sphere}
\end{figure}

\FloatBarrier

\paragraph{Example: the projective line is a circle.}
For $n=1$, the map $(x,y)\mapsto(x^2-y^2,2xy)$ from $S^1$ onto $S^1$
has exactly antipodal fibers (it doubles the angle). Its continuous factor
is a bijection from a compact space to a Hausdorff space, proving
$\mathbb{RP}^1\cong S^1$.

\subsection{Affine patches and their coordinates}
A line with $x_i\ne0$ meets the affine plane $x_i=1$ in exactly one point.
The remaining coordinates of that intersection give the desired chart.
\begin{definition}[Standard affine patches]
For $i=0,\ldots,n$, define
\[
U_i=\{[x_0:\cdots:x_n]:x_i\ne0\},\qquad
\varphi_i:U_i\to\R^n,\quad \varphi_i([x])=(x_j/x_i)_{j\ne i},
\]
with the retained indices in increasing order.
\end{definition}
\begin{proposition}[Affine charts]
The $U_i$ are open, cover $\mathbb{RP}^n$, and each $\varphi_i$ is a
homeomorphism onto $\R^n$.
\end{proposition}
\begin{proof}
\emph{Domains and representative independence.}
The preimage $\pi^{-1}(U_i)=\{x:x_i\ne0\}$ is open in
$\R^{n+1}\setminus\{0\}$, so $U_i$ is open. Every nonzero vector has
a nonzero coordinate, proving the cover assertion. Rescaling multiplies
both numerator and denominator by the same nonzero scalar, so the ratios
and their domains are independent of the representative.

\emph{Continuity by descent.} On $\{x_i\ne0\}$ define the continuous
ratio map
\[
\gamma_i(x)=(x_j/x_i)_{j\ne i}.
\]
 The restriction of $\pi$
over $U_i$ is quotient: a preimage open in this restricted source is open
in the full source, since $\pi^{-1}(U_i)$ is open; the quotient criterion
then applies. The universal property gives
$\widetilde{\gamma_i}=\varphi_i$, proving continuity.

\emph{Continuous inverse.} For $u=(u_j)_{j\ne i}\in\R^n$, let $s_i(u)$
insert $1$ in position $i$, keeping all other entries in increasing index
order. This is a continuous map into $\R^{n+1}\setminus\{0\}$, and
\[
\varphi_i^{-1}=\pi\circ s_i.
\]
Indeed, taking ratios of $s_i(u)$ returns $u$, while inserting $1$ into
the ratios of $x$ returns
\[
x/x_i,
\]
 representing the same line as $x$.
This verifies both inverse identities and continuity of the inverse.
\end{proof}

\begin{example}[The first overlap]
For $n\geq1$, if $\varphi_0([x])=(u_1,\ldots,u_n)$, then $U_0\cap U_1$
corresponds to $u_1\ne0$. Dividing instead by $x_1$ gives
\[
\varphi_1\circ\varphi_0^{-1}(u_1,\ldots,u_n)
=\left(\frac1{u_1},\frac{u_2}{u_1},\ldots,\frac{u_n}{u_1}\right).
\]
For $n=1$, only the first entry remains.
\end{example}
\paragraph{General transitions.}
For $i\ne j$, the domain of $\varphi_j\circ\varphi_i^{-1}$ is
$\{u:u_j\ne0\}$. Its output has indices $k\ne j$, in increasing order:
\[
v_i=1/u_j,\qquad v_k=u_k/u_j\quad(k\ne i,j).
\]
These rational maps are smooth where their denominators are nonzero;
the inverse exchanges $i$ and $j$.

\begin{proposition}[Hausdorffness from the antipodal metric]
The space $\mathbb{RP}^n$ is Hausdorff.
\end{proposition}
\begin{proof}
Apply \cref{prop:isometric-quotient-metric} to $Y=S^n$ and the two
isometries $\operatorname{id}$ and $-\operatorname{id}$. The finiteness
hypothesis is automatic, and $\|x-(-x)\|=2$. The metric
\[
d([x],[y])=\min\{\|x-y\|,\|x+y\|\}\qquad(x,y\in S^n)
\]
induces the antipodal quotient topology, already identified with the
scaling quotient. Every metric space is Hausdorff.
\end{proof}

\paragraph{Second countability, compactness, and smooth structure.}
Pull back the countable Euclidean basis of \cref{ex:rational-basis} in
each of the $n+1$ open charts.
Their finite union is countable and is a basis: intersect any open
neighborhood with a chart containing its point and choose a basis member
there. Thus the space is second countable. It is compact because the
continuous map
\[
S^n\to S^n/{\pm1}\cong\mathbb{RP}^n
\]
 is onto and $S^n$
is compact. The smooth transitions give an atlas; extending it maximally
proves that $\mathbb{RP}^n$ is a compact smooth $n$-manifold without boundary.
For $n=0$ the space is a point with a chart into $\R^0$.


\Needspace{24\baselineskip}
\begin{example}[The disk model of $\mathbb{RP}^2$]
Every antipodal pair on $S^2$ has a representative in the closed upper
hemisphere; the representative is unique in its interior, while equatorial
points remain paired antipodally. Projection of the hemisphere onto its
first two coordinates is a homeomorphism onto $D^2$, with inverse
$u\mapsto(u,\sqrt{1-\|u\|^2})$. The induced map from the disk quotient
to $\mathbb{RP}^2$ is a continuous bijection. Compactness of the disk
quotient and the Hausdorff property just proved make it a homeomorphism.
Thus $\mathbb{RP}^2$ identifies only antipodal boundary pairs, whereas

\[
D^2/\partial D^2\cong S^2
\]
 collapses the entire boundary to one point.
\end{example}
\begin{figure}[htbp]
\centering
\begin{tikzpicture}[font=\small,>=Stealth]
\foreach \cx in {0,4.6}{\fill[gray!6] (\cx,0) circle(1);}
\draw[blue!75!black,very thick] (0,0) circle(1);
\foreach \a in {0,45,...,315}{\fill[blue!75!black] ({cos(\a)},{sin(\a)}) circle(.055);}
\node[above,align=center] at (0,1.2) {$S^2$: collapse\\the entire boundary};
\node[below,align=center] at (0,-1.2) {all marked points\\form one class};
\draw[very thick] (4.6,0) circle(1);
\foreach \a in {30,210}{\fill[blue!75!black] ({4.6+cos(\a)},{sin(\a)}) circle(.065);}
\foreach \a in {110,290}{\draw[red!75!black,thick,fill=white] ({4.6+cos(\a)-.065},{sin(\a)-.065}) rectangle ++(.13,.13);}
\draw[dotted,blue!75!black] ({4.6+cos(30)},{sin(30)})--({4.6-cos(30)},{-sin(30)});
\draw[dotted,red!75!black] ({4.6+cos(110)},{sin(110)})--({4.6-cos(110)},{-sin(110)});
\node[above,align=center] at (4.6,1.2) {$\mathbb{RP}^2$: pair only\\antipodal boundary points};
\node[below,align=center] at (4.6,-1.2) {each marker shape\\is a separate pair};
\end{tikzpicture}
\caption{Two different disk quotients. For the sphere, all of $\partial D^2$
collapses together. For the projective plane, $u\sim-u$ only on the
boundary; interior points stay distinct. Dotted chords indicate pairings,
not curves collapsed in the disk.}
\label{fig:projective-disk}
\end{figure}

\FloatBarrier

\Needspace{10\baselineskip}
\section*{Exercises}
\addcontentsline{toc}{section}{Exercises}
\begin{hexercise}{app-e-1}
Show that $g(t)=(\cos 2\pi t,\sin 2\pi t)$ descends to $\R/\Z$.
Explain separately why the descended map is well defined, continuous,
and a homeomorphism.
\end{hexercise}
\begin{exercise}
Descend $(x,t)\mapsto(e(x),t)$ from the rectangle to $Q_1$ and prove
$Q_1\cong S^1\times[-1,1]$. Identify its two boundary components.
\end{exercise}
\begin{hexercise}{app-e-3}
For $\sigma=-1$, write the seam chart centered at
\[
[0,1/2]
\]
 using a disk
of radius
\[
1/8.
\]
 State its two finite-rectangle domains and its inverse.
Repeat at $[0,1]$, using half-disk coordinates.
\end{hexercise}
\begin{exercise}
Trace the M\"obius boundary from $(0,1)$ through both horizontal edges.
Write a parametrization with period $2$ and prove that its classes have
no smaller positive period.
\end{exercise}
\begin{exercise}
For $\mathbb{RP}^2$, compute $\varphi_1\varphi_0^{-1}(u_1,u_2)$,
state its domain, and check its inverse explicitly.
\end{exercise}
\begin{hexercise}{app-e-6}
Verify that $(x,y)\mapsto(x^2-y^2,2xy)$ on $S^1$ has exactly antipodal
fibers and descends to a homeomorphism $\mathbb{RP}^1\to S^1$.
\end{hexercise}
